\section{Implementation}\label{sec:implementation}

We implemented the method in PyTorch without any reference code.
Each stage is one script, and a separate module handles prototype fitting, log-densities and sampling.

\paragraph{Networks.}
We follow the architectures in the paper's supplementary material.
For Digits, $f_s$ is a modified LeNet with BatchNorm and 100-dimensional features.
For Office-31 and VisDA, $f_s$ is an ImageNet-pretrained ResNet-50~\citep{he2016resnet}.
A head of three fully connected layers (1024, 256 and 256 units) with ELU activations and BatchNorm follows it and produces 256-dimensional features.
The classifiers $g_s$ and $g_t$ are two-layer MLPs.
The projections $f_e$ and $f_d$ have two fully connected layers for Digits and three for Office and VisDA.

\subsection{Choices the paper leaves open}\label{sec:choices}
\Cref{tab:choices} lists the settings the paper does not give, together with our choices.
We tuned the Digits epoch counts on M$\to$U only and then kept them fixed for the other Digits tasks.

\begin{table}[t]
  \centering
  \small
  \caption{Settings not specified in the paper, and our choices.}
  \label{tab:choices}
  \begin{tabularx}{\columnwidth}{@{}>{\raggedright\arraybackslash}p{2.0cm}>{\raggedright\arraybackslash}X@{}}
    \toprule
    Setting & Our choice and reason \\
    \midrule
    Prototype score $\ell^c$ & $\log\mathcal{P}^c(\uu)$, because raw densities in 100 to 256 dimensions underflow. \\
    Covariance & $\bm{\Sigma}\leftarrow 0.9\,\bm{\Sigma} + (0.1+10^{-4})\,\tfrac{\operatorname{tr}\bm{\Sigma}}{D}\,\mathbf{I}$. An Office class has approximately 30 to 90 images in 256 dimensions, so its raw covariance is singular. \\
    Batch size & 64 for Digits, 32 for Office and VisDA. \\
    Source epochs & 10 (Digits), 50 (Office), 5 (VisDA). \\
    Adaptation epochs & 30 (Digits), 60 (Office), 10 (VisDA). \\
    Adaptation epoch length & Five times the number of target batches, so each of the five losses sees approximately one pass over the target data. \\
    Projection pretraining & 2{,}000 iterations. \\
    Data splits & Digits adapts on the target training split and tests on the target test split. Office adapts and tests on the whole target domain, minus the labelled images. \\
    Warm-up & Updates both $f_s$ and $g_s$. The algorithm lists only $f_s$, but $g_s$ also needs training. \\
    Private guides & One-shot images pass through $f_t$ in the same batch as target samples, so BatchNorm never sees a batch of only one image per class. \\
    \bottomrule
  \end{tabularx}
\end{table}

\subsection{Our changes to the method}\label{sec:changes}
When we ran the method as written, shared and private classes became mixed during adaptation (\cref{sec:exp-vanilla}).
We made four changes, each aimed at one cause, and \cref{sec:exp-improvements} reports their effect.
A configuration preset (\texttt{preset=vanilla}) switches all four off.

\paragraph{1. Target BatchNorm statistics.}
Because $f_t$ starts as a copy of $f_s$, it also inherits the BatchNorm running statistics of the source domain.
The first pseudo-labels are computed with these statistics, which creates a mismatch.
The private guides come from target images, while the shared guides come from source means.
As a result, many shared target samples lie closer to a private guide than to their own shared guide.
Before adaptation, we therefore re-estimate the BatchNorm statistics of $f_t$ on the unlabelled target set, a technique known as AdaBN~\citep{li2018adabn}.
This fits the design of the method, since $f_t$ is meant to be specific to the target.

\paragraph{2. Projections that start as the exact identity.}
The paper pretrains $f_e$ and $f_d$ as plain MLPs to approximate the identity, but some error remains after pretraining.
This error moves shared-class features across the decision boundaries of $g_s$ before adaptation even starts.
The supplementary material also puts an activation on the last layer, which distorts the negative values that BatchNorm features often take.
ELU clips them below $-1$ and LeakyReLU scales them by $0.01$, so the identity is hard to reach.
We instead use $f_e(\uu)=\uu+h(\uu)$, where $h$ is an MLP whose linear last layer is initialised to zero, and we give $f_d$ the same form.
Both projections are then exactly the identity at the start, and pretraining begins from zero error.

\paragraph{3. $\mathcal{L}_{a1}$ on the confident set.}
As written, $\mathcal{L}_{a1}$ \eqref{eq:la1} uses every target sample, each labelled by its nearest guide.
Early in adaptation many of these labels are wrong, so $\mathcal{L}_{a1}$ trains $f_t$ to fit its own errors.
This pulls shared samples into private clusters, and the effect grows with every epoch.
We therefore draw the $\mathcal{L}_{a1}$ batches from the confident set $\Bt$ instead.
The label is still the nearest guide computed online, and $\Pt^c$ is still fitted on all target samples.

\paragraph{4. Frozen ResNet-50 for Office and VisDA.}
The paper only states that the ResNet-50 is pretrained on ImageNet.
Fine-tuning the whole network on approximately 900 source images of 10 classes reaches 100\% source accuracy.
However, it discards the features that separate the unseen target classes.
We therefore freeze the ResNet-50 weights in both stages, while its BatchNorm statistics still adapt.
The fully connected head, the classifiers, $f_e$ and $f_d$ train as before.

\subsection{Inconsistencies in the paper}\label{sec:inconsistencies}

\paragraph{Office private classes.}
The supplementary material lists the private classes as those ranked 21 to 31, which gives 11 classes.
The main results table, however, uses $|\Ct|=20$, which implies 10 private classes.
We therefore use classes 21 to 30 for the setting $|\Cs|=10, |\Ct|=20$, and classes 21 to 31 for $|\Cs|=20, |\Ct|=31$.

\paragraph{The $3\sigma$ rule in high dimensions.}
The Mahalanobis distance measures how far a point lies from a Gaussian's mean, in units of its covariance.
Under $\mathcal{N}(\bm{\mu},\bm{\Sigma})$, the point $\bm{\mu}+3\sqrt{\lambda}\,\bm{\nu}$ in \eqref{eq:neg} has a squared Mahalanobis distance of exactly $9$.
The rule therefore keeps a candidate whenever its Mahalanobis distance to the most likely class exceeds~$3$.
For a sample from a $D$-dimensional Gaussian, however, the squared distance follows a $\chi^2_D$ distribution with mean $D$.
With $D=100$ or $256$, almost every sample exceeds $9$, including real samples of the class itself.
The rule thus keeps almost every candidate as a negative.

\paragraph{Typo in the negative-sampling algorithm.}
Line~9 of Algorithm~1 in the supplementary material writes $\mathcal{P}(\uu_s)$ where $\mathcal{P}(\uu_n)$ is meant.

\subsection{Engineering}\label{sec:engineering}
Each run saves a checkpoint every epoch and every 20 minutes, and it resumes automatically after an interruption.
Cluster jobs have a 6-hour limit, so they requeue themselves before reaching it.
Unit tests check prototype fitting, log-densities and sampling, and a short end-to-end test runs both stages on synthetic data.
Diagnostic scripts count how many shared samples are predicted as private and the reverse.
They also track pseudo-label accuracy per epoch, measure nearest-guide accuracy on fixed features, and plot confusion matrices.
